\documentclass[paper=a4, fontsize=11pt]{scrartcl}

\usepackage[T1]{fontenc}
\usepackage{fourier}
\usepackage[english]{babel}
\usepackage{amsmath,amsfonts,amsthm}

\usepackage{sectsty}
\allsectionsfont{\centering \normalfont\scshape}

\usepackage{fancyhdr}
\pagestyle{fancyplain}
\fancyhead{}
\fancyfoot[L]{}
\fancyfoot[C]{}
\fancyfoot[R]{\thepage}
\renewcommand{\headrulewidth}{0pt}
\renewcommand{\footrulewidth}{0pt}
\setlength{\headheight}{13.6pt}

\setlength\parindent{0pt}

% Title helpers
\newcommand{\horrule}[1]{\rule{\linewidth}{#1}}

\title{
\normalfont \normalsize
\textsc{Artificial Intelligence - 4,906,1.00 - Sheet 1 Spring 2026} \\[10pt]
\horrule{0.5pt} \\[0.4cm]
\huge Solutions \\[-2pt]
\horrule{2pt} \\[0.5cm]
}

\author{
\textbf{Group 12}\\
Valeria Pointet\\
Keanu Belo da Silva
}

\date{\normalsize\today}

\begin{document}
\maketitle

% -------------------- 1 --------------------
\section{Turing Test and Intelligence}

\subsection*{Question 1.1}
The Turing Test (1950) is a text-only conversation test.
A human interrogator chats with a human and a machine.
Both try to act like a human.
The interrogator tries to tell which one is the machine and which one is the human.
The purpose is to test whether a machine can appear human-like in conversation.

\subsection*{Question 1.2}
A key shortcoming is that the test checks only behavior in text.
It does not prove real understanding.
The Chinese Room thought experiment explains this:
a system can follow rules to produce correct answers without understanding the meaning.
So a machine could pass the test by manipulating symbols, not by truly ``thinking''.

\subsection*{Question 1.3}
We could make the test harder by requiring more than short text chat.
For example, add real-world interaction or multi-modal inputs (images, audio), and run it longer.
Then pure imitation is harder, and the system must stay consistent and competent over time.

\subsection*{Question 1.4}
Correct answer: Evaluating whether an AI system is more intelligent than a human.

The Turing Test is not a ranking test.
It only checks whether the system can be mistaken for a human in conversation.

% -------------------- 2 --------------------
\section{Strong vs.\ Weak AI}

\subsection*{Question 2.1}
Weak AI (Narrow AI) is AI in the engineering sense.
It is built for specific tasks and appears intelligent.
Example: a spam filter or a face recognition system.

Strong AI (Artificial General Intelligence) is AI in the philosophical sense.
It would be self-conscious and have human-level general intelligence.
Example: a truly general, self-conscious machine (mostly hypothetical today).

\subsection*{Question 2.2}
ChatGPT is Weak AI.
It can appear intelligent in language tasks, but it is not considered self-conscious.
So it fits Narrow AI, not Strong AI.

% -------------------- 3 --------------------
\section{Agent Environments}

\subsection*{Question 3.1}
PEAS stands for: Performance measure, Environment, Actuators, Sensors.

For a package delivery drone:

\textbf{P (Performance measure):}
safe delivery, correct location, fast delivery, low energy use, no crashes, legal compliance.

\textbf{E (Environment):}
buildings, trees, people, birds, other drones, weather (wind/rain), no-fly zones, GPS noise.

\textbf{A (Actuators):}
rotors for movement, altitude control, turning, package release mechanism, emergency landing.

\textbf{S (Sensors):}
camera, GPS, IMU (accelerometer/gyroscope), altitude sensor, proximity sensors (e.g.\ lidar),
battery sensors.

\subsection*{Question 3.2}
Observability: \textbf{partially observable}.
Sensors do not give the complete state (e.g.\ hidden obstacles, sensor noise).

Determinism: \textbf{stochastic}.
Wind and other agents create uncertainty, so the next state is not fully determined.

Episodic vs.\ sequential: \textbf{sequential}.
Each action affects later states (route, safety, battery).

Static vs.\ dynamic: \textbf{dynamic}.
The environment can change while the drone is deciding (moving objects, changing weather).

Discrete vs.\ continuous: \textbf{continuous}.
Position, velocity, time, and wind change continuously.

Single-agent vs.\ multi-agent: \textbf{multi-agent}.
Other drones, people, and traffic act independently.

% -------------------- 4 --------------------
\section{Agent Types}

\subsection*{Question 4.1}
A rational agent chooses the action that is expected to maximize its performance measure,
given what it has perceived and what it knows.
Rational does not mean omniscient.
An omniscient agent would know the complete true state of the environment.
A rational agent can still have missing information, so rational $\neq$ omniscient.

\subsection*{Question 4.2}
An agent that plays chess needs planning/search.
The following agent types would not be able to handle chess well:

\textbf{Simple reflex-based agent:}
it uses condition-action rules from the current percept only.
It does not plan ahead, so it cannot reliably handle deep tactics.

\textbf{Model-based reflex agent:}
it keeps an internal state, but it still selects actions like a reflex agent.
Without search/lookahead, it also cannot handle chess well.

Chess is typically tackled by agents that use \textbf{search} (goal-based / utility-based)
and often also \textbf{learning}.

\end{document}
